\section{Modelos de negocio de la conducción autónoma: quién es dueño de los datos y de los clientes}

Después de exponer la trayectoria personal de entrenamiento en Waymo y Momenta, hay que pasar al nivel del modelo de negocio. Al igual que en la robótica, en la conducción autónoma la ruta tecnológica nunca es una elección puramente técnica: está ligada al punto de acceso a los datos, a la relación con los clientes y a la bolsa de utilidades.

En la entrevista, Gao Jiyang (高继扬) divide los modelos de negocio de la conducción autónoma en varias categorías: el robotaxi al estilo de Waymo, la venta de vehículos más suscripción de software de los fabricantes de automóviles, el proveedor al estilo de Momenta y la plataforma de utilidades del vehículo completo, similar a la de Huawei. Esta clasificación ayuda a entender por qué se bifurcan las rutas tecnológicas: quién es dueño del vehículo, quién es dueño del usuario, quién es dueño de los datos y quién asume la responsabilidad del servicio influye directamente en la arquitectura técnica y en las capacidades de la organización.

\begin{figure}[H]
\centering
\includegraphics[width=0.94\textwidth]{autonomous-driving-business-models.png}
\caption{Los cuatro modelos de negocio de la conducción autónoma: Robotaxi, suscripción del fabricante de automóviles, proveedor y plataforma de utilidades del vehículo completo.}
\end{figure}

\begin{knowledgebox}{Lectura de la figura: el modelo de negocio determina el ciclo de los datos}
El Robotaxi opera su propia flota: los datos y el servicio están en sus manos, pero su escalamiento es lento. La suscripción del fabricante de automóviles posee los vehículos y los usuarios, y obtiene los datos directamente. El proveedor necesita los proyectos de los fabricantes para llegar a la producción en serie, y tanto el valor para el cliente como el retorno de los datos dependen de la colaboración. La plataforma de utilidades del vehículo completo, por su parte, redefine la bolsa de utilidades del vehículo mediante la marca, los canales de venta, la conducción inteligente y la cabina inteligente.
\end{knowledgebox}

\subsection{Por qué esto es instructivo para la robótica}

La clasificación anterior proviene de la conducción autónoma, pero esta sección busca trasladarla: la robótica no cuenta con un mercado ya establecido como el automotriz ni con datos de flota de origen natural. Precisamente porque carece de estas condiciones previas, la experiencia de la conducción autónoma solo puede ofrecer un marco, no respuestas directas.

La inteligencia corpórea todavía no tiene un lado de la demanda ni una escala de flota tan maduros como los del automóvil. Lo que la conducción autónoma enseña a la robótica no es copiar el modelo de robotaxi o el de proveedor, sino ver con claridad que el ciclo cerrado de datos debe estar ligado al valor para el cliente. Si un robot no tiene escenarios reales ni unidades entregadas, no hay datos continuos; sin datos, es difícil entrenar un modelo fundacional de acciones; sin la capacidad del modelo, no se puede aumentar el valor para el cliente.

\begin{warningbox}{La experiencia de la conducción autónoma no puede copiarse mecánicamente}
El automóvil ya tiene un mercado enorme: los vehículos se venden. El cuerpo del robot de propósito general todavía no cuenta con una demanda igual de madura. La robótica debe crear al mismo tiempo la demanda de hardware, el valor de los escenarios y el ciclo cerrado de datos, lo cual es más difícil que «trasladar la ruta de la conducción autónoma».
\end{warningbox}

\subsection{Resumen de la sección}

El problema de fondo de los modelos de negocio de la conducción autónoma es: de dónde provienen los datos, cómo se forma el valor para el cliente y cómo itera el sistema. La inteligencia corpórea tampoco puede eludir estas preguntas; la diferencia es que su demanda, su hardware y sus escenarios están más dispersos.
